\documentclass{amsart}
% For dealing with references we use the comment environment
\usepackage{verbatim}
\newenvironment{reference}{\comment}{\endcomment}
%\newenvironment{reference}{}{}
\newenvironment{slogan}{\comment}{\endcomment}
\newenvironment{history}{\comment}{\endcomment}

% For commutative diagrams we use Xy-pic
\usepackage[all]{xy}

% We use 2cell for 2-commutative diagrams.
\xyoption{2cell}
\UseAllTwocells

% We use multicol for the list of chapters between chapters
\usepackage{multicol}

% This is generally recommended for better output
\usepackage{lmodern}
\usepackage[T1]{fontenc}

% For cross-file-references

% Package for hypertext links:
\usepackage{hyperref}

% For any local file, say "hello.tex" you want to link to please

% Theorem environments.
%
\theoremstyle{plain}
\newtheorem{theorem}[subsection]{Theorem}
\newtheorem{proposition}[subsection]{Proposition}
\newtheorem{lemma}[subsection]{Lemma}

\theoremstyle{definition}
\newtheorem{definition}[subsection]{Definition}
\newtheorem{example}[subsection]{Example}
\newtheorem{exercise}[subsection]{Exercise}
\newtheorem{situation}[subsection]{Situation}

\theoremstyle{remark}
\newtheorem{remark}[subsection]{Remark}
\newtheorem{remarks}[subsection]{Remarks}

\numberwithin{equation}{subsection}

% Macros
%
\def\lim{\mathop{\mathrm{lim}}\nolimits}
\def\colim{\mathop{\mathrm{colim}}\nolimits}
\def\Spec{\mathop{\mathrm{Spec}}}
\def\Hom{\mathop{\mathrm{Hom}}\nolimits}
\def\Ext{\mathop{\mathrm{Ext}}\nolimits}
\def\SheafHom{\mathop{\mathcal{H}\!\mathit{om}}\nolimits}
\def\SheafExt{\mathop{\mathcal{E}\!\mathit{xt}}\nolimits}
\def\Sch{\mathit{Sch}}
\def\Mor{\mathop{\mathrm{Mor}}\nolimits}
\def\Ob{\mathop{\mathrm{Ob}}\nolimits}
\def\Sh{\mathop{\mathit{Sh}}\nolimits}
\def\NL{\mathop{N\!L}\nolimits}
\def\CH{\mathop{\mathrm{CH}}\nolimits}
\def\proetale{{pro\text{-}\acute{e}tale}}
\def\etale{{\acute{e}tale}}
\def\QCoh{\mathit{QCoh}}
\def\Ker{\mathop{\mathrm{Ker}}}
\def\Im{\mathop{\mathrm{Im}}}
\def\Coker{\mathop{\mathrm{Coker}}}
\def\Coim{\mathop{\mathrm{Coim}}}

% Boxtimes
%
\DeclareMathSymbol{\boxtimes}{\mathbin}{AMSa}{"02}

%
% Macros for moduli stacks/spaces
%
\def\QCohstack{\mathcal{QC}\!\mathit{oh}}
\def\Cohstack{\mathcal{C}\!\mathit{oh}}
\def\Spacesstack{\mathcal{S}\!\mathit{paces}}
\def\Quotfunctor{\mathrm{Quot}}
\def\Hilbfunctor{\mathrm{Hilb}}
\def\Curvesstack{\mathcal{C}\!\mathit{urves}}
\def\Polarizedstack{\mathcal{P}\!\mathit{olarized}}
\def\Complexesstack{\mathcal{C}\!\mathit{omplexes}}
% \Pic is the operator that assigns to X its picard group, usage \Pic(X)
% \Picardstack_{X/B} denotes the Picard stack of X over B
% \Picardfunctor_{X/B} denotes the Picard functor of X over B
\def\Pic{\mathop{\mathrm{Pic}}\nolimits}
\def\Picardstack{\mathcal{P}\!\mathit{ic}}
\def\Picardfunctor{\mathrm{Pic}}
\def\Deformationcategory{\mathcal{D}\!\mathit{ef}}

\setlength{\emergencystretch}{3em}
\begin{document}
\title[Stacks Project, Tag 0B10]{460 --- Cycle-level projection formula for flat proper maps of nonsingular varieties}
\maketitle
\noindent Copyright (C) 2005--2025 Johan de Jong. Stacks Project authors. Source: \href{https://stacks.math.columbia.edu/tag/0B10}{Tag 0B10}.

\noindent Editorial difficulty note (not author text). Difficulty H2: The proof must establish proper intersection on the target and match multiplicities before passing to Chow groups. It separates positive-dimensional generic fibres from generically finite supports, uses properness to obtain local finiteness near each intersection component, and compares Tor lengths through the derived projection formula at those stalks..

\noindent Editorial provenance note (not author text). History: excerpt from revision \texttt{a0444\allowbreak 6e57e\allowbreak c1fbc\allowbreak 252a8\allowbreak 71afc\allowbreak ec775\allowbreak 2fb28\allowbreak 07b14}, intersection.tex, lines 2042--2143. Reference presentation was adapted; original-author context and core support blocks were retained or added. The mathematical wording is unchanged. Licensed under GNU FDL 1.2 or later; no invariant sections or cover texts. See ../licenses/COPYING. Context blocks reproduce author definitions and supporting results; they are not additional counted problems.

\section{Conventions}
\label{section-conventions}

\noindent
We fix an algebraically closed ground field $\mathbf{C}$ of any
characteristic. All schemes and varieties are over $\mathbf{C}$ and all
morphisms are over $\mathbf{C}$. A variety $X$ is
{\it nonsingular} if $X$ is a regular scheme (see
Properties, Definition \ref{properties-definition-regular}).
In our case this means that the morphism $X \to \Spec(\mathbf{C})$
is smooth (see
Varieties, Lemma \href{https://stacks.math.columbia.edu/tag/038X}{lemma geometrically regular smooth (Tag 038X)}).




\section{Intersection product using Tor formula}
\label{section-intersection-product}

\noindent
Let $X$ be a nonsingular variety. Let
$\alpha = \sum n_i [W_i]$ be an $r$-cycle and
$\beta = \sum_j m_j [V_j]$ be an $s$-cycle on $X$.
Assume that $\alpha$ and $\beta$ intersect properly, see
Definition \ref{definition-proper-intersection}.
In this case we define
$$
\alpha \cdot \beta = \sum\nolimits_{i,j} n_i m_j W_i \cdot V_j.
$$
where $W_i \cdot V_j$ is as defined in Section \ref{section-tor-formula}.
If $\beta = [V]$ where $V$ is a closed subvariety of dimension $s$,
then we sometimes write $\alpha \cdot \beta = \alpha \cdot V$.



\begin{definition}
\label{properties-definition-regular}
Let $X$ be a scheme. We say $X$ is {\it regular}, or {\it nonsingular} if
for every $x \in X$ there exists an affine open neighbourhood
$U \subset X$ of $x$ such that the ring $\mathcal{O}_X(U)$ is
Noetherian and regular.
\end{definition}

\begin{definition}
\label{definition-proper-intersection}
Let $X$ be a nonsingular variety.
\begin{enumerate}
\item Let $W,V \subset X$ be closed subvarieties with
$\dim(W) = s$ and $\dim(V) = r$. We say that $W$ and $V$
{\it intersect properly} if $\dim(V \cap W) \leq r + s - \dim(X)$.
\item Let $\alpha = \sum n_i [W_i]$ be an $s$-cycle,
and $\beta = \sum_j m_j [V_j]$ be an $r$-cycle on $X$. We say
that $\alpha$ and $\beta$ {\it intersect properly} if
$W_i$ and $V_j$ intersect properly for all $i$ and $j$.
\end{enumerate}
\end{definition}

\section{Intersection multiplicities using Tor formula}
\label{section-tor-formula}

\noindent
A basic fact we will use frequently is that given sheaves of
modules $\mathcal{F}$, $\mathcal{G}$ on a ringed space $(X, \mathcal{O}_X)$
and a point $x \in X$ we have
$$
\text{Tor}_p^{\mathcal{O}_X}(\mathcal{F}, \mathcal{G})_x =
\text{Tor}_p^{\mathcal{O}_{X, x}}(\mathcal{F}_x, \mathcal{G}_x)
$$
as $\mathcal{O}_{X, x}$-modules. This can be seen in several ways
from our construction of derived tensor products in
Cohomology, Section \href{https://stacks.math.columbia.edu/tag/06Y7}{section flat (Tag 06Y7)}, for example it follows from
Cohomology, Lemma \href{https://stacks.math.columbia.edu/tag/06YB}{lemma check K flat stalks (Tag 06YB)}.
Moreover, if $X$ is a scheme and $\mathcal{F}$ and $\mathcal{G}$
are quasi-coherent, then the modules
$\text{Tor}_p^{\mathcal{O}_X}(\mathcal{F}, \mathcal{G})$ are
quasi-coherent too, see
Derived Categories of Schemes, Lemma
\href{https://stacks.math.columbia.edu/tag/08DX}{lemma quasi coherence tensor product (Tag 08DX)}.
More important for our purposes is the following result.



\begin{lemma}
\label{lemma-projection-formula-flat}
\begin{reference}
See \href{https://stacks.math.columbia.edu/bibliography}{Bibliography: Serre\_algebre\_locale (Chapter V, C), Section 7, formula (10))}
for a more general formula.
\end{reference}
Let $f : X \to Y$ be a flat proper morphism of nonsingular varieties.
Set $e = \dim(X) - \dim(Y)$. Let $\alpha$ be an $r$-cycle on $X$ and let
$\beta$ be a $s$-cycle on $Y$. Assume that $\alpha$ and $f^*(\beta)$ intersect
properly. Then $f_*(\alpha)$ and $\beta$ intersect properly and
$$
f_*(\alpha) \cdot \beta = f_*( \alpha \cdot f^*\beta)
$$
\end{lemma}

\begin{proof}
By linearity we reduce to the case where $\alpha = [V]$ and
$\beta = [W]$ for some closed subvariety $V \subset X$ and
$W \subset Y$ of dimension $r$ and $s$. Then $f^{-1}(W)$ has
pure dimension $s + e$. We assume the cycles
$[V]$ and $f^*[W]$ intersect properly. We will use without
further mention the fact that $V \cap f^{-1}(W) \to f(V) \cap W$
is surjective.

\medskip\noindent
Let $a$ be the dimension of the generic fibre of $V \to f(V)$.
If $a > 0$, then $f_*[V] = 0$. In particular $f_*\alpha$ and $\beta$
intersect properly. To finish this case we have to show that
$f_*([V] \cdot f^*[W]) = 0$. However, since every fibre of
$V \to f(V)$ has dimension $\geq a$ (see
Morphisms, Lemma \href{https://stacks.math.columbia.edu/tag/02FZ}{lemma openness bounded dimension fibres (Tag 02FZ)})
we conclude that every irreducible component $Z$ of $V \cap f^{-1}(W)$
has fibres of dimension $\geq a$ over $f(Z)$. This certainly
implies what we want.

\medskip\noindent
Assume that $V \to f(V)$ is generically finite. Let $Z \subset f(V) \cap W$
be an irreducible component. Let $Z_i \subset V \cap f^{-1}(W)$,
$i = 1, \ldots, t$ be the irreducible components of $V \cap f^{-1}(W)$
dominating $Z$. By assumption each $Z_i$ has dimension
$r + s + e - \dim(X) = r + s - \dim(Y)$. Hence
$\dim(Z) \leq r + s - \dim(Y)$. Thus we see that $f(V)$ and $W$
intersect properly, $\dim(Z) = r + s - \dim(Y)$, and each
$Z_i \to Z$ is generically finite. In particular, it follows that
$V \to f(V)$ has finite fibre over the generic point $\xi$ of $Z$.
Thus $V \to Y$ is finite in an open neighbourhood of $\xi$, see
Cohomology of Schemes, Lemma
\href{https://stacks.math.columbia.edu/tag/02OH}{lemma proper finite fibre finite in neighbourhood (Tag 02OH)}.
Using a very general projection formula for derived tensor products, we get
$$
Rf_*(\mathcal{O}_V \otimes_{\mathcal{O}_X}^\mathbf{L} Lf^*\mathcal{O}_W) =
Rf_*\mathcal{O}_V \otimes_{\mathcal{O}_Y}^\mathbf{L} \mathcal{O}_W
$$
see Derived Categories of Schemes, Lemma
\href{https://stacks.math.columbia.edu/tag/08EU}{lemma cohomology base change (Tag 08EU)}.
Since $f$ is flat, we see that $Lf^*\mathcal{O}_W = f^*\mathcal{O}_W$.
Since $f|_V$ is finite in an open neighbourhood of $\xi$ we have
$$
(Rf_*\mathcal{F})_\xi = (f_*\mathcal{F})_\xi
$$
for any coherent sheaf on $X$ whose support is contained in $V$
(see Cohomology of Schemes, Lemma
\href{https://stacks.math.columbia.edu/tag/02OE}{lemma higher direct images zero finite fibre (Tag 02OE)}). Thus
we conclude that
\begin{equation}
\label{equation-stalks}
\left(
f_*\text{Tor}_i^{\mathcal{O}_X}(\mathcal{O}_V, f^*\mathcal{O}_W)
\right)_\xi =
\left(\text{Tor}_i^{\mathcal{O}_Y}(f_*\mathcal{O}_V, \mathcal{O}_W)\right)_\xi
\end{equation}
for all $i$. Since $f^*[W] = [f^*\mathcal{O}_W]_{s + e}$ by
Lemma \href{https://stacks.math.columbia.edu/tag/0AZF}{lemma pullback (Tag 0AZF)} we have
$$
[V] \cdot f^*[W] =
\sum (-1)^i
[\text{Tor}_i^{\mathcal{O}_X}(\mathcal{O}_V,
f^*\mathcal{O}_W)]_{r + s - \dim(Y)}
$$
by Lemma \href{https://stacks.math.columbia.edu/tag/0B0W}{lemma tor sheaf (Tag 0B0W)}. Applying
Lemma \href{https://stacks.math.columbia.edu/tag/0AZD}{lemma push coherent (Tag 0AZD)}
we find
$$
f_*([V] \cdot f^*[W]) =
\sum (-1)^i
[f_*\text{Tor}_i^{\mathcal{O}_X}(\mathcal{O}_V,
f^*\mathcal{O}_W)]_{r + s - \dim(Y)}
$$
Since $f_*[V] = [f_*\mathcal{O}_V]_r$ by
Lemma \href{https://stacks.math.columbia.edu/tag/0AZD}{lemma push coherent (Tag 0AZD)} we have
$$
[f_*V] \cdot [W] =
\sum (-1)^i
[\text{Tor}_i^{\mathcal{O}_X}(f_*\mathcal{O}_V,
\mathcal{O}_W)]_{r + s - \dim(Y)}
$$
again by Lemma \href{https://stacks.math.columbia.edu/tag/0B0W}{lemma tor sheaf (Tag 0B0W)}.
Comparing the formula for $f_*([V] \cdot f^*[W])$ with
the formula for $f_*[V] \cdot [W]$ and looking at the
coefficient of $Z$ by taking lengths of stalks at $\xi$, we see that
(\ref{equation-stalks}) finishes the proof.
\end{proof}
\end{document}
